% Section 7 -- Limitations. Ported from paper/limitations.md (L1-L6) with the
% material the four-page limit cut restored. Ordered by how much damage each
% objection does if we do not raise it first. No literal numbers: every figure
% is a macro from numbers.tex.

\section{Limitations}
\label{sec:limitations}

Three objections are load-bearing, and we state them first:
the corpora target different mathematics (L4), the instrument correction is
measured on one side and transferred to the other (L1), and everything rests on
one machine corpus (L6).

\paragraph{L1. The correction transfers a human rate to a larger population.}
Machine-side miss rates are imputed rather than measured, from per-class rates
measured on human proofs and standardized onto the machine tactic mix.
Arithmetic-tactic proofs are \machineArithMixPct{}\% of the machine corpus and
\humanArithMixPct{}\% of mathlib, so most of the machine corpus is standardized
against a stratum that is a fraction of a percent of mathlib. That rate is
near-censused and precise; the exposure is that the two populations'
\texttt{nlinarith} may not be the same object. In mathlib it is typically one
step closing a residual goal; in the machine corpus it is frequently the whole
proof. If the extractor's blind spot scales with how much of a proof
the automation closes rather than with its presence, the transfer
under-corrects, leaving machine proofs looking narrower than they are, which
biases toward our own hypothesis. Flat efficiency across human tactic classes
(Section~\ref{sec:reference-line}) constrains the instrument's behaviour on
mathlib proofs only; nothing measured here establishes comparable recovery on
machine proofs, whose dominant stratum has almost no human counterpart. We
therefore treat the correction as a lower bound and require the gap to clear a
multiple of it, not merely exceed it. Elaboration-level extraction on the machine corpus would settle it. We did not
do it because it means reproducing the publishers' extraction pipeline on a
second corpus (building the premise tool against the same toolchain and
wrapping each standalone proof as a compilable module), and because those
proofs were verified against their releasers' mathlib rather than our pinned
commit, so a fraction would fail on renames and deprecations, giving version
skew rather than a miss rate.

\paragraph{Instrument limitations.} Four further entries are stated in
Appendix~\ref{app:limitations}: L2, the asymmetries left in the accounting and
why every one of them narrows the reported gap; L3, how much of the explicit set source-visible extraction reaches; L3d,
that Table~\ref{tab:retention}'s quantities are matcher-based; and L3e, that
our reading of the flag is not validated against elaborated terms.

\paragraph{L4. The corpora target different mathematics.} Machine corpora
target competition statements and mathlib targets library development, so part
of L1's tactic-mix difference belongs to the problem sets, not the provers.
Section~\ref{sec:size-controls} narrows this by matching library
\emph{area} without removing it, because area is not \emph{genre}. A mathlib
theorem is general-purpose infrastructure; a competition problem is a
self-contained puzzle. As an exploratory, not pre-registered check, mathlib's own IMO archive gives
\genreCheckRatio{}$\times$, the smallest gap of any design we ran. It is a
lower bound (most of those declarations cite local helper lemmas outside the
graph's build target), it is not extraction-audited, and it is not
regenerable from the released artifact (L6). We report it because we measured
it, and publishing only the designs that give larger gaps would be the
objectionable choice. We therefore claim a difference
between these two released corpora under the restrictions named, not a general
property of how provers use mathlib when proving comparable theorems.

\paragraph{L6. Scope, and what is not regenerable.} One machine corpus,
\machineProofs{} proofs. These are comparisons of two distributions and not
causal claims. We write ``machine proofs concentrate on'', never ``provers are
biased toward''. Two families of figures carried from the prior version are \emph{not}
reproducible from the released artifact: the per-premise occurrence shares of
Section~\ref{sec:concentration} and the genre-check figures in L4. Both were
computed ad hoc, neither table was
committed, and no released script regenerates them. Every other number comes from a committed result file, bar
Section~\ref{sec:dag-layer}'s graph-audit figures, which any reader recomputes
from the released graph in one pass. We flag them rather than recompute them,
because recomputing after the fact is how a frozen scope stops being one.
